% Self-contained normalization and raw/normal-form comparison.
% Adapted with explicit attribution from Research/full_normalization and
% incoming ipra_witness_discharge_2026-11-11/report.tex, normalization appendix.
% New explanatory addition: the administrative-let proof for the strict
% normalizer, including evaluation of an argument erased by the update.

\section{Normalization and a direct bridge to raw Choice-DI}
\label{sec:normalization-bridge}
\label{sec:canonical-bridge}

This section reconstructs the normalization and raw-comparison arguments
of the repository's full-normalization package~\cite{normalization} and the
incoming witness-discharge report~\cite{witness}. The
compound-envelope repair is inherited from those sources. The direct
comparison below connects the present deterministic Name calculus to the
raw Choice-DI foundations included later, without taking the incoming
witness-discharge equivalence as a black box. The explicit
administrative-\texttt{let} argument also accounts for the selected
normalizer's strict evaluation of a recursive value that its update may
subsequently discard.

All assertions about normalization in this section are external theorems
about finite syntax. They are not new arithmetic axioms or an
object-language uniform evaluator.

\subsection{Strong normalization of scalar computation}

On raw terms $x,0,St,I(n,a;u.h)$ let $\longrightarrow$ be the contextual
closure, including under update binders, of
\[
 I(0,a;u.h)\longrightarrow a,\qquad
 I(Sn,a;u.h)\longrightarrow h[I(n,a;u.h)/u].
 \tag{Raw computation}\label{eq:nb-reduction}
\]
Bound variables are freshened before substitution. These reductions are
stable under capture-avoiding substitution. Call a finite term strongly
normalizing if it has no infinite reduction sequence. For such a term $t$,
write $\nu(t)$ for the largest length of a reduction beginning at $t$ and
$|t|$ for its syntax-tree size. The maximum exists because the reduction
tree is finitely branching and has no infinite path.

\begin{lemma}[Closure under iteration]
\label{lem:nb-iterator}
Suppose $n,a$ are strongly normalizing and $h[v/u]$ is strongly normalizing
for every strongly normalizing $v$. Then $I(n,a;u.h)$ is strongly
normalizing.
\end{lemma}
\begin{proof}
Taking $v=u$ shows that $h$ is strongly normalizing. Induct
lexicographically on
\[
            (\nu(n),\ |n|,\ \nu(a)+\nu(h)).
\]
It suffices to show that every one-step reduct of the iterator is strongly
normalizing. A contraction inside $n$ decreases the first coordinate.
A contraction inside $a$ or $h$ decreases the last. If $h\to h'$, then
$h'[v/u]$ is a reduct of $h[v/u]$ for every $v$, so the universal body
hypothesis persists. The zero root reduct is $a$. For a successor root
$n=Sp$, we have $\nu(p)=\nu(n)$ and $|p|<|n|$. The induction hypothesis
makes $I(p,a;u.h)$ strongly normalizing, and the hypothesis about $h$
then applies to the successor root reduct
$h[I(p,a;u.h)/u]$.
\end{proof}

\begin{theorem}[Strong normalization]
\label{thm:nb-sn}
Replacing the free variables of any fixed finite raw term by strongly
normalizing terms produces a strongly normalizing term. In particular,
every finite raw term is strongly normalizing.
\end{theorem}
\begin{proof}
Induct on the fixed raw term, using simultaneous capture-free substitutions.
Variables, zero, and successor are immediate. At an iterator, first rename
its bound state variable away from the finitely many substituting terms.
The structural induction hypotheses give strongly normalizing counter and
initializer instances. The induction hypothesis for the strictly smaller
update term applies to every additional strongly normalizing assignment
for its state variable. It gives exactly the universal premise of
Lemma~\ref{lem:nb-iterator}. The identity substitution yields the last
assertion.
\end{proof}

\subsection{Confluence and normalization--substitution coherence}

\begin{lemma}[Confluence]
\label{lem:nb-confluence}
The reduction~\eqref{eq:nb-reduction} is confluent. Every term has a unique
normal form $N(t)$, up to alpha-renaming, and literally with de Bruijn
indices.
\end{lemma}
\begin{proof}
First check local confluence. Disjoint redex contractions commute, and
the zero and successor root patterns do not overlap. If an inner redex
lies inside a root contraction, contract each of its residual copies on
the other path. A zero contraction returns its initializer and erases
the entire update; erased residuals require no step. A successor
contraction may copy its recursive iterator into
several update positions. The update is finite, so only finitely many
copies arise, and substitution stability gives a common reduct after
contracting them. An inner contraction of the update itself is handled
both in the outer substituted update and in the retained recursive
iterator. These are all nested overlaps.

Strong normalization and local confluence imply confluence: induct on
the reduction height of a common source to join the first steps of two
finite reductions and then their remaining tails. Thus two normal forms
reachable from a term are equal up to bound-variable naming.
\end{proof}

A term is normal exactly when its children are normal and no iterator
counter is headed by $0$ or by $S$. Formula normalization changes only
numerical terms; the equality and $D$ constructors are preserved.

\begin{lemma}[Coherence]
\label{lem:nb-coherence}
For every finite raw term $t$ and capture-free substitution $\sigma$,
\[
 N(t[\sigma])=
 N\bigl(N(t)[x\mapsto N(\sigma x)]\bigr).
 \tag{Coherence}\label{eq:nb-coherence}
\]
The same equality holds for formulas. Normalized substitution consequently
obeys identity, composition, renaming, and the usual binder-lifting laws.
\end{lemma}
\begin{proof}
The reduction $t\to^*N(t)$ survives raw substitution. Within the finitely
many substituted copies, reduce each $\sigma x$ to $N(\sigma x)$.
The resulting terms are convertible, and confluence makes their normal
forms identical. Applying this observation twice and using composition
of raw substitutions proves
\[
 N\bigl(N(t[\sigma])[\tau]\bigr)
   =N\bigl(t[x\mapsto N((\sigma x)[\tau])]\bigr).
\]
Identity, renaming, and lifting are special cases. The formula assertion
follows structurally because formula constructors do not reduce.
\end{proof}

\paragraph{The strict implementation also terminates.}
The selected traversal normalizes the counter, initializer, and update
first. At a successor counter it normalizes the predecessor iterate
before inserting it into the update. If the update ignores its state,
that recursive value would have been erased by the direct raw root
contraction, so this strict evaluation deserves an explicit justification.

Introduce an administrative term $\mathsf{let}\ u=a\ \mathsf{in}\ h$,
with its usual capture-free reduction to $h[a/u]$, and replace the
successor root rule in the auxiliary syntax by
\[
 I(Sn,a;u.h)\longrightarrow
       \mathsf{let}\ v=I(n,a;u.h)\ \mathsf{in}\ h[v/u],
 \tag{Strict expansion}
\]
where $v$ is fresh. Reduce under all constructors and let components.
The zero rule is unchanged. There are no new normal forms: every let
has a root contraction, and its contraction composes a strict expansion
to the original successor computation.

For this auxiliary reduction, a let is strongly normalizing whenever
its value is strongly normalizing and its body sends every strongly
normalizing substitution to a strongly normalizing term. Induction on
the sum of the component reduction heights proves this: internal
contractions decrease the sum, and the root contraction is covered by
the universal body hypothesis. The proof of
Lemma~\ref{lem:nb-iterator} now uses this let-closure fact at its successor
root, after obtaining the predecessor iterate by the same lexicographic
induction. Structural induction, also including the let constructor,
therefore proves strong normalization of the auxiliary syntax.

The strict traversal is an evaluation strategy in this auxiliary
calculus: after strict expansion, it first normalizes the let value and
then contracts the let. Thus its recursive computational work cannot
continue indefinitely, even in the erased-value case. Between those
contractions it performs finite tree traversals and index renamings.
Forgetting administrative lets preserves conversion in the original
calculus; its returned normal term is therefore the unique $N(t)$.
The same argument and~\eqref{eq:nb-coherence} identify hereditary
substitution with normalized capture-free substitution. No uniform
termination theorem inside iPRA has been invoked.

\subsection{The raw reference calculus and its harmless clause convention}

Let $\mathcal C_{\mathrm{raw}}$ be raw Choice-DI over the terms of
\eqref{eq:nb-reduction} and formulas $s=t,D_{x,y}B$. It has the two raw
computation equations, formula transport, clause Use, ordinary paired
elimination, and the two separate introductions
\[
 \frac{\Gamma,B(t,0)\vdash B(t,Sz)}{\Gamma\vdash D_{x,y}B},
 \qquad
 \frac{\Gamma,E(t,z)\vdash E(t,Sz)}{\Gamma\vdash D_{x,y}E},
 \tag{Raw introductions}\label{eq:nb-raw-intro}
\]
where the second matrix is a literal equation and all eigenvariable
conditions are as in the inherited Choice-DI foundations.

For comparison we allow clauses with zero or one antecedent and retain
$\mathcal R=[\,;y.y=y]$. This is only an interface convention. In the
inherited single-antecedent raw calculus, replace an unconditional clause
$[\,;y.B]$ by $[H;y.B]$ and supply the theorem $H$ to each use.
Every use of $\mathcal R$ is replaced by derived raw reflexivity.
Conversely the optional-premise interface admits all single-antecedent
clauses. The inherited computation theorem used as a neutral antecedent
there and $H$ are both theorems, so those ordinary-assumption encodings
are interchangeable by inserting the corresponding neutral proof.
No new arithmetical or logical axiom is added by this convention.

The included Choice-DI development derives equation induction with
arbitrary passive contexts, numerical cases, all ordinary equality
reasoning, and binder congruence. Its logical and scalar arithmetic
bridges identify its ordinary realization with full first-order iPRA.
We use those results, rather than an assertion only about standard
natural-number models.

\subsection{Raw provability of normalization, including update binders}

Transport in a numerical parameter does not permit a free variable in
an equality proof to become captured by an iterator update. The needed
update congruence follows from literal-equation induction.

\begin{lemma}[Update extensionality]
\label{lem:nb-update}
If $u$ is fresh for $\Gamma$ and
$\Gamma\vdash_{\mathcal C_{\mathrm{raw}}}h(u)=k(u)$, then, after
freshening binders away from the passive data,
\[
 \Gamma\vdash_{\mathcal C_{\mathrm{raw}}}
     I(n,a;u.h)=I(n,a;u.k)
\]
for every counter $n$ and initializer $a$.
\end{lemma}
\begin{proof}
Apply derived literal-equation induction to
$E(z):=(I(z,a;u.h)=I(z,a;u.k))$, with $z$ fresh for the retained data.
At zero both sides compute to $a$. In the successor step, write the
left and right predecessor iterates as $L$ and $R$. Computation gives
$h(L)$ and $k(R)$. Capture-free numerical substitution in the supplied
pointwise proof gives $h(L)=k(L)$; the induction hypothesis $L=R$
and ordinary term transport give $k(L)=k(R)$. Transitivity and the
two successor computations prove $E(Sz)$. Induction gives $E(n)$.
Only the literal equation $E$ is inducted on.
\end{proof}

\begin{lemma}[Internal normalization equivalence]
\label{lem:nb-internal}
For every fixed finite raw term $t$,
\[
              \vdash_{\mathcal C_{\mathrm{raw}}}t=N(t).
\]
For every fixed raw formula $A$, the formulas $A$ and $N(A)$ are
interderivable, uniformly under additional retained contexts.
\end{lemma}
\begin{proof}
Each root contraction is one raw computation equation. Ordinary
transport handles a contraction inside a successor, counter, or
initializer. For a contraction beneath an update binder, use
Lemma~\ref{lem:nb-update}, repeating it for nested update contexts.
Thus every step in a finite normalization sequence is internally a
provable equality. Transitivity gives $t=N(t)$.

For formulas, induct externally on formula depth. At an equation,
normalize its two terms using the proved term equalities. At
$D_{x,y}B$, open the source formula at a fresh witness, introduce the
target at that same witness by case introduction, convert the new
zero base to the old one by the induction hypothesis, use the old
clause at a fresh successor, and convert the result back. Reverse the
conversions for the reverse implication. Substitution may create fresh
redexes in both an instance of $B$ and the corresponding instance of
$N(B)$: apply the smaller-depth induction hypothesis to both instances
and use coherence to identify their common normal form. These are uniform
derivations under every passive context. The construction uses
derived binder congruence; it never captures a free variable by an
ordinary equality transport.
\end{proof}

\paragraph{Replacing a retained clause.}
If antecedents $A,A'$ are uniformly interderivable and bodies $B(y),B'(y)$
are uniformly interderivable, then the clauses $[A;y.B]$ and
$[A';y.B']$ are interchangeable in derivations. Replace each use of the
old clause by conversion to the new antecedent, use of the new clause,
and conversion of its specialized conclusion. Induction on proof trees,
with fresh local eigenparameters, justifies this under further retained
clauses. The same statement holds for unconditional clauses, with no
antecedent conversion. In particular, a raw context and its normalized
context have the same consequences after the corresponding formula
conversions. This supplies the context-conversion step explicitly,
rather than presuming that ordinary-assumption cut alone rewrites a
conditional-universal clause.

\subsection{Repairing the omitted case introduction at witness zero}

Write $\mathcal A$ for the deterministic three-family calculus specified
in Section~\ref{sec:deterministic-name}. Its previously derived ordinary
elimination and numerical cases require no equality transport. For unused
$u$, put $q_v F(v)=D_{u,v}F(v)$. For an equation $E=(s=t)$ define
\[
 G_vE=(C(v,0,s)=C(v,0,t)),\qquad L(E)=q_vG_vE,
\]
with all terms normalized and all new variables fresh. Direct computation
gives $G_0E\equiv H$ and $G_{Sw}E\equiv E$.

\begin{lemma}[Compound envelope]
\label{lem:nb-envelope}
In $\mathcal A$, an equation $E$ and its compound envelope $L(E)$ are
interderivable under every retained context, without equality transport.
\end{lemma}
\begin{proof}
Under $E$, introduce the unused witness of $q_vG_vE$ at $1$. Its
successor premise is literally $E$, so this is case introduction.
Conversely, open $q_vG_vE$ by derived ordinary elimination. Its clause
has antecedent $H$. Prove $H$ by the fixed reflexivity clause and use
the opened clause at $1$ to obtain $E$.
\end{proof}

\begin{lemma}[Unrestricted case introduction]
\label{lem:nb-all-cases}
The case introduction in~\eqref{eq:nb-raw-intro}, with normalized
substitution throughout, is derivable in $\mathcal A$ for every witness
and every matrix.
\end{lemma}
\begin{proof}
Only a literal-equation matrix $E(x,y)$ at the literal witness zero
is missing from the primitive case mode. Suppose
$\Gamma,E(0,0)\vdash E(0,Sz)$ is supplied. Put
$A(x,y)=L(E(x,y))$ with suitably fresh envelope binders. Under
$A(0,0)$, unwrap its envelope, insert the supplied derivation by
ordinary-assumption cut, and wrap the result. This gives $A(0,Sz)$.
Since $A$ has a $D$ root, primitive Intro at zero now selects cases,
and yields $D_{x,y}A(x,y)$.

Open it by ordinary elimination at a fresh $c$. Introduce
$D_{x,y}E(x,y)$ at witness $c$. Its syntax is a parameter rather than
literal zero, so Intro again selects cases. Under $E(c,0)$, wrap the
equation to $A(c,0)$, use the opened clause at $Sz$, and unwrap
$A(c,Sz)$ to $E(c,Sz)$. Closing Intro and ordinary elimination gives
the desired formula. No claim that $c\ne0$ is used.
\end{proof}

\begin{lemma}[Equation progress at every witness]
\label{lem:nb-all-progress}
The equation-progress introduction in~\eqref{eq:nb-raw-intro}, with
normalized substitution, is derivable in $\mathcal A$ for every witness.
\end{lemma}
\begin{proof}
Given $\Gamma,E(t,z)\vdash E(t,Sz)$, use a fresh unused witness
coordinate $w$. Primitive Intro at zero on the equation matrix
$E(t,y)$ yields $D_{w,y}E(t,y)$ by progress. Under the prospective
outer case assumption $E(t,0)$, open this formula by ordinary
elimination and use its clause at $Sz$. This provides the successor
premise for unrestricted case introduction of $D_{x,y}E(x,y)$
at witness $t$, available by Lemma~\ref{lem:nb-all-cases}.
The unused coordinate prevents the opened witness from replacing $t$.
\end{proof}

\subsection{An auxiliary normal-form calculus}

Let $\mathcal C_{\mathrm{nf}}^+$ use canonical terms and formulas,
normalized substitution, optional-premise clauses and the fixed
reflexivity clause. Its rules are Use, ordinary paired elimination,
full formula transport, and the separate unrestricted case and
literal-equation progress introductions. There is no separate
computation axiom, since both computation equations normalize to
reflexivity.

\begin{proposition}[Canonical Choice-DI and the new calculus]
\label{prop:nb-canonical-choice}
The calculi $\mathcal C_{\mathrm{nf}}^+$ and $\mathcal A$ derive the
same canonical sequents.
\end{proposition}
\begin{proof}
Translate $\mathcal C_{\mathrm{nf}}^+$ to $\mathcal A$ by using its
derived ordinary elimination, term and formula transport, and
Lemmas~\ref{lem:nb-all-cases}--\ref{lem:nb-all-progress}. Use and the
fixed reflexivity clause are shared.

Conversely each primitive Intro is one of the two available choices.
For Name, ordinary elimination opens its witness clause. When the
deterministic guard succeeds, Use at $1$, supplied with reflexivity
of zero, proves the selected equality; otherwise reflexivity proves
$b=b$. Full formula transport gives the exported normalized formula,
whose stipulated freedom from $b$ permits ordinary elimination to
close. All translations preserve the entire canonical end sequent.
\end{proof}

\subsection{Both proof translations and both composites}

Let $T$ normalize all numerical terms in raw formulas and clauses.
Let $U$ regard canonical syntax as raw syntax. These are operations on
finite descriptions, not extra operations on the numerical domain.

\begin{theorem}[Direct raw/canonical comparison]
\label{thm:direct-raw-canonical}
For every raw sequent and every canonical sequent,
\begin{align*}
 \Gamma\vdash_{\mathcal C_{\mathrm{raw}}}A
  &\quad\Longleftrightarrow\quad
        T\Gamma\vdash_{\mathcal A}T A,\\
 \Delta\vdash_{\mathcal A}B
  &\quad\Longleftrightarrow\quad
        U\Delta\vdash_{\mathcal C_{\mathrm{raw}}}U B.
\end{align*}
Moreover $TU(B)=B$ literally, and $A$ and $UT(A)$ are uniformly
interderivable in the raw calculus. The corresponding clause-context
composites are interchangeable in every derivation.
\end{theorem}
\begin{proof}
First translate a raw derivation into $\mathcal C_{\mathrm{nf}}^+$
by induction on derivations. The two computation equations have
identical normal forms on their two sides, so become reflexivity.
Coherence~\eqref{eq:nb-coherence} makes the other normalized
premises and conclusions instances of the corresponding canonical
rules. Equation-root status is preserved, and normalization cannot
introduce a new free eigenparameter. Crucially, the auxiliary calculus
retains both introduction choices, so this translation never assumes
that a nonzero-syntax raw witness stays nonzero after normalization.
Proposition~\ref{prop:nb-canonical-choice} then translates to
$\mathcal A$.

In the reverse direction first translate $\mathcal A$ to
$\mathcal C_{\mathrm{nf}}^+$. Regard its canonical formulas as raw.
For a raw rule whose substitution instance is not already normal,
Lemma~\ref{lem:nb-internal} converts the normalized premises to raw
instances, the raw rule is applied, and its conclusion is converted
back. The retained clause of an elimination may have a normalized
zero antecedent instead of the literal raw substitution instance;
the explicit clause-replacement construction above converts that
context. The same reasoning handles Intro's temporary ordinary
assumption. Raw reflexivity supplies the fixed reflexivity-clause
uses. Thus induction on canonical derivations gives a raw proof
with the stated end sequent.

On canonical syntax $TU$ is literal identity. On raw syntax $UT$
is formula normalization, whose internal equivalence and clause
interchangeability were proved above. Composing the two proof
translations with these equivalences gives both reverse implications
in the displayed statements as well as the composites.
\end{proof}

\begin{corollary}[Full first-order bi-interpretation]
\label{cor:nb-full-ipra}
Composing the direct comparison with the included Choice-DI logical and
scalar arithmetic bridges gives a bi-interpretation of the present
calculus with full first-order intuitionistic primitive recursive
arithmetic. Both numerical domains are whole, both equalities are
actual equality, and both numerical composite isomorphisms have graph
$x=y$. Both complete formula composites are internally provable.
The theories are equiconsistent.
\end{corollary}
\begin{proof}
The inherited bridge reconstructs all intuitionistic logical
connectives and both quantifiers, quantifier-free induction, and the
fixed-description primitive-recursive function symbols, with both
formula round trips. Theorem~\ref{thm:direct-raw-canonical} supplies
the remaining change of term presentation. Its operations commute
with substitution modulo internally provable computation equalities,
so it preserves those interpretations. Numerical variables normalize
to themselves: the comparison changes finite term descriptions,
not the numerical domain or equality. Hence the inherited numerical
identity graphs remain $x=y$.

The common designated contradiction $0=S0$ is already normal and
is preserved in both directions. The full formula comparisons
therefore imply equiconsistency, as well as the stronger
bi-interpretability statement.
\end{proof}

\subsection{Uniform normalization is external}

The full first-order iPRA convention here has a symbol for each fixed
primitive-recursive definition and the corresponding defining
equations, with quantifier-free induction in intuitionistic logic.
It is neither the equation-only fragment nor a theory with unrestricted
formula induction. A fixed scalar term compiles to one such
primitive-recursive definition; no numerical variable ranges over all
definitions inside the object theory.

Uniform normalization over arbitrary finite scalar-term codes is total
recursive but is not primitive recursive. Otherwise, by coding a unary
scalar term, substituting a numeral, normalizing, and extracting the
resulting numeral, one would obtain a primitive-recursive universal
evaluator $V(e,n)$ for all primitive-recursive functions. Invalid codes
could be assigned zero. The diagonal function $V(e,e)+1$ would itself
have a scalar-term code, contradicting universality. This uses the
included scalar compilation theorem, not an assumed uniform evaluator.

The size of a complete checker must count its normalizer, and the
unbounded interface has no primitive-recursive running-time bound.
Explicit finite fuel or reduction traces give a different,
primitive-recursive bounded certificate interface for the same
theorems. They do not add an arithmetic assertion of a uniform bound
on the necessary fuel. The metatheorems here are written proofs; the
finite executable checks are implementation evidence rather than a
new Lean or Rocq verification of them.
